\chapter*{Заключение}
\addcontentsline{toc}{chapter}{Заключение}
\label{ch:conclusion}

В ходе выполнения лабораторной работы~\cfgLabNumber{} были решены
четыре независимых учебных задания на базовые приёмы HTML и CSS в
онлайн-редакторе \textbf{Codepen}. Все четыре стартовых pen-а
форкнуты под собственным аккаунтом, изменения сохранены; параллельно
поддерживаются локальные копии решений в каталоге
\texttt{lab12/codepen-solutions/}.

Освоенные понятия и приёмы:
\begin{itemize}
  \item семантическая разметка текста --- заголовки шести уровней,
        абзацы, цитаты \texttt{<blockquote>}, гиперссылки
        \texttt{<a>} и изображения \texttt{<img>} с обязательным
        атрибутом \texttt{alt} (блок~1.1);
  \item списки трёх типов --- маркированные \texttt{<ul>},
        нумерованные \texttt{<ol>}, описательные \texttt{<dl>};
        правила сортировки данных в нумерованных списках по
        смыслу (блок~1.2);
  \item таблицы с группировкой ячеек --- атрибуты \texttt{colspan}
        и \texttt{rowspan}; типовой класс ошибок <<развалившейся>>
        таблицы, когда число колонок в разных строках не совпадает,
        и приёмы её починки (блок~1.2);
  \item CSS-селекторы разной специфичности --- селектор по тегу,
        по классу, селектор потомка через пробел, прямой потомок
        через \texttt{>}; стратегии выбора селектора в условиях,
        когда нужно затронуть только часть одинаковых элементов
        (блок~2.1);
  \item типографические свойства --- \texttt{color},
        \texttt{font-size}, \texttt{font-family},
        \texttt{font-style}, \texttt{font-weight},
        \texttt{text-decoration}, \texttt{list-style-type}
        с двумя уровнями маркеров для вложенных списков,
        \texttt{background-image} с относительным URL (блок~2.2).
\end{itemize}

Полученные навыки замыкают переход от дизайн-составляющей курса
(лабораторные работы~10 и~11, где макеты собирались в Figma) к
непосредственной реализации интерфейса кодом. Концепции лабораторных
работ~10--11 имеют прямые аналоги в HTML/CSS: фреймы Figma ---
\texttt{<div>} и \texttt{<section>}; Auto-Layout --- свойство
\texttt{display:~flex}; Constraints (привязки) --- свойство
\texttt{position} и медиа-запросы; Text~styles --- классы CSS;
Color~styles --- CSS-переменные. Это делает плавным дальнейший
переход к практической вёрстке многостраничного сайта.

\endinput
